% !TEX root = ../meu-trabalho.tex
% -----------------------------------------------------------------------------
% Introdução
% -----------------------------------------------------------------------------

\chapter{Introdução}
\label{chap_introducao}
\section{Contextualização}
\label{sec_contextualizacao}

``A administração financeira pode ser vista como uma forma de economia aplicada, que se baseia amplamente em conceitos teóricos econômicos'' \citeonline[p.~4]{gitman2004principles}. Com essa citação, se introduz o presente artigo abordando a definição de sua área de estudo, a Administração Financeira. Área que se caracteriza por dinamismo, utilizando de novas tecnologias em busca de maximização de resultados.

Nesse sentido, \citeonline[p.~15]{paiva2014redes}, diz que a possibilidade de prever preços futuros em ativos financeiros, contraria a teoria financeira dominante, a Hipótese de Mercado Eficiente (HME). Esta, denota ineficácia na previsão de preços do mercado se utilizando de informações passadas. A afirmativa é justificada pelo pressuposto de movimento browniano da HME, concebido por características randômicas e independentes. Entretanto, Paiva pontua que há estudos que defendem que o mercado acionário, na verdade, não se firma como aleatório e as séries financeiras têm memória de longo-prazo. De tal maneira que, caso seja possível identificar os padrões de comportamento do mercado, se tornaria plausível uma abordagem preditiva.

Sendo assim, diversos estudos demonstraram obter sucesso em previsões de preços de ativos, fazendo a utilização de ferramentas de aprendizado de máquina. Em sua definição clássica, \citeonline[p.~2]{mitchell1997machine} define aprendizado de máquina como a capacidade de um programa de computador aprender com experiência E, com respeito a uma determinada classe de tarefas T e em seguida, aferir a partir de uma métrica de desempenho P se seu desempenho em T melhora com experiência E.

\citeonline[p.~331]{aquart2022machine}, empregou e analisou vários modelos de Aprendizado de máquina para prever e negociar movimentos das 100 maiores criptomoedas concluiu que todos os modelos utilizados fazem previsões estatisticamente viáveis.

\citeonline[p.~852]{wei2023forecasting} descobriram que, ao combinar a previsão de retornos do Bitcoin com uma estratégia de alavancagem híbrida que utiliza volatilidade e sentimento narrativo, que segundo o texto são informações de sentimento extraídas de fontes textuais formais e públicas, como publicações, notícias e artigos de negócios globais, que influenciam o mercado de criptomoedas, em contraste com dados brutos de mídias sociais ou indicadores técnicos. Todas as estratégias de alavancagem melhoraram significativamente o desempenho da negociação, sendo que a estratégia híbrida superou as demais.

\citeonline{belcastro2023enhancing} afirma que a partir da previsão de preços de criptomoedas através da integração de dados de mídias sociais e de mercado, alcançou um ganho médio geral de 194\% sem taxas de transação e de 117\% considerando as taxas.

\section{Problema de Pesquisa}
\label{sec_problema_de_pesquisa}

Tendo em vista o cenário citado acima, a presente pesquisa busca responder: quais as contribuições da aprendizagem de máquina na gestão de carteiras de criptoativos?

\section{Objetivos de pesquisa}
\label{sec_objetivos_de_pesquisa}

\subsection{Objetivo geral}
\label{subsec_objetivo_geral}

Diante de tal pergunta, a pesquisa há de examinar as contribuições do uso de aprendizado de máquina para a gestão de carteiras de criptoativos.

\subsection{Objetivos específicos}
\label{subsec_objetivos_especificos}

Firmando se a esse objetivo, delinearam-se os seguintes objetivos específicos para viabilizar a acurácia das determinações da pesquisa:

\begin{alineas}
    \item Analisar os resultados obtidos a partir do modelo de diversificação de portifólio de Markowitz em conjunto com o aprendizado de máquina.
    \item Comparar estratégias de investimento tradicionais e as com base em modelos de aprendizado de máquina.
    \item Avaliar se há consistência entre os resultados apresentados em diferentes pesquisas.
\end{alineas}

\section{Justificativa}
\label{sec_justificativa}

Em um contexto de ascensão do termo ``inteligência artificial'', é difícil saber factualmente o que pode ser promissor na utilização de novas tecnologias computacionais baseadas em interações não lineares. Segundo \citeonline[p.~140]{molnar2019interpretable}, são interações não lineares,

\begin{citacao}
Quando características interagem entre si em um modelo de previsão, a previsão não pode ser expressa como a soma dos efeitos das características, porque o efeito de uma característica depende do valor da outra característica.
\end{citacao}

Por isso, a presente pesquisa tem por objetivo analisar as contribuições do uso de aprendizado de máquina para a gestão de carteiras de criptoativos. Tal objetivo se justifica, pois é capaz de servir de base para a criação de ferramentas, para a análise de viabilidade de investimentos e diversificação de carteira. Sendo assim, espera-se contribuir com a tomada de decisão de gestores, tanto bancários quanto de departamentos financeiros: ``O \textit{Machine Learning} financeiro corresponde a um conjunto de procedimentos para estimar um modelo estatístico e utilizar esse modelo para tomar decisões.''~\cite[p.~10]{kelly2023financial}.

A pesquisa aprofunda sobre áreas de utilização de aprendizado de máquina, e coloca em pauta um grande debate da economia, a Hipótese de Mercado eficiente (HME). Tal teoria, pressupõe que os resultados do mercado financeiro não permitem ganhos acima do normal, pois, segundo \citeonline{fama1970efficient}, na HME os preços dos ativos refletem integralmente toda a informação disponível, teoria que passa por recorrentes contestações. Além disso, trabalhos científicos anteriores, já abordavam a necessidade do aprofundamento da pesquisa em Machine learning e suas utilizações, como em \citeonline{sarker2021machine} e \citeonline{clara2024survey}. Ademais, a pesquisa espera ser benéfica para o contexto social, contribuindo com a criação ferramentas que tragam maior racionalidade e métodos mais robustos para a tomada de decisões no mercado financeiro. De tal maneira, evitando que haja a formação de bolhas, que segundo \citeonline{scherbina2013asset}, ocorrem quando preços de ativos sofrem desvios em relação ao valor fundamental dos ativos. A pesquisa espera também que seus achados tenham contribuição para as ODS, em particular a ODS 9 --- Indústria, inovação e infraestrutura, a partir do debate sobre novas tecnologias e seu uso de forma responsável no desenvolvimento de infraestruturas de dados.
